%! Author = shoupu_wan
%! Date = 2/2/20

\section{Case Study for Data Structure Design}\label{sec:case-study-data-structure-design}

This problem is an optimization problem concerning the capital gain of a company.
Starting with an initial capital \code{cap}, the company has a to-do list, projects out of which
the company has targeted to finish \code{k} of them.
There are two lists of equal size, \code{capitals} and \code{profits}, which specify the required
capital and net profit for each corresponding project.
The question is how to select the projects and in what order to finish them so as to maximize the
final capital gain.

\subsection{Na\"{i}ve Solution}\label{subsec:max-cap-naive}
First of all and to help reader warm up on this problem, the following greedy algorithm would
not work and is a false start.
\begin{enumerate}
    \item Heapify the projects with the smallest capital and largest profit on top
    \item heappop projects
    \item If enough capital to finish the project, collect the profits
    \item If not, return
\end{enumerate}

What's wrong with this?
It fails to handle the sufficient capital but limited project case.

So really this problem deals with two criteria at once---attainability and profitability.
Attainability criteria has to do with the subset of projects attainable with the current capital.
Profitability criteria choose the most profitable projects among those attainable.
Because each project can only be done once, one need to update the subset of attainable projects
each time a project is completed.
Based on these observation, we can lay out the first algorithm
\begin{alisting}
    \caption{Na\"{i}ve algorithm for the `Max Capital' problem.\label{max-capital-naive-algo}}
    \begin{enumerate}
        \item Sort projects by `(capital, profit descending)'
        \item Find the most profitable one among the attainable projects
        \item Collect the profit and delete the project from projects to-do list
        \item Repeat \code{k} times or until attainable projects run out
    \end{enumerate}
\end{alisting}

A \python implementation of this algorithm is listed in~\fullref{lst:max-capital-brute-force}.
The part removing a finished project from to-be-considered list is costly.
So despite the simplicity of this solution, the runtime for this algorithm is $O(N^2)$.
\begin{clisting}
    \lstinputlisting[label={lst:max-capital-brute-force},
    caption={Max capital gain.
    A naive implementation}]
    {max-capital-brute-force.py}
\end{clisting}

At this time it is unclear how to beat that runtime.
But looking back at the two criteria---attainability and profitability---we may decouple these
two requirements and meet each of them with separate data structures.
Therefore we may process the projects in a streamlined fashion with two stages---with the first
stage procuring attainable projects and the second extracting the most
profitable among the attainable.
Also note that the total capital is a non-decreasing sequence with each finished project.
This means that if a project was attainable before will continue to be thereafter.
This is important to our above streamlined design---projects can only advance along the path and
never going back or cycle back to the beginning.

the heap may be chosen for both stage.

\begin{clisting}
    \lstinputlisting[label={lst:max-capital-two-hop},
    caption={Max capital gain.
    A two-data-structure
    implementation}]{max-capital-two-hops.py}
\end{clisting}
Unlike the na\"{i}ve solution, this data structure design separates the concerns for attainability
from profitability.
